% !TEX root = ../main.tex
\section{Impact/Financial Benefits}
\label{sec:impact}

% Revised 2026-10-11 from the owner's hand review of the printed report
% (handwritten/2026-10-11_main.md, notes 20.x and 21.x): two sentences struck from
% the scope flag; the emissions subsection moved to Section "Emissions and Who
% Pays" in 06_results_discussion.tex.

\plantodo{S9, S12 (S16 bullet added this session)}{Quantify the impact of the reference architecture and its experiment. Translate the conflict analysis and the experiment's results into concrete operational terms wherever possible (e.g., delay-minutes, fuel burn, CO2, instructions issued) and into US-dollar terms where a credible conversion exists (e.g., published cost-per-delay-minute figures from the literature review, Section~\ref{sec:litreview}). Also discuss the ethical context and any unintended consequences of the recommendations (equity-minded engineering: whose costs get externalized, and onto whom -- this connects directly to the objectives whose costs are externalized onto another stakeholder, Section~\ref{sec:objectives-results}, and to Section~\ref{sec:emissions-impact}).}

\textbf{Scope flag:} the 503 template expects a quantified financial-benefit statement (the template's sample statements are dollar-savings claims from industry DMAIC-style projects). This capstone is an architecture/SoS project, not a cost-reduction project with a ``before'' baseline in dollars, so a literal dollar-savings figure may not be available or meaningful. Recommend framing impact primarily in operational/systems terms (delay, fuel, emissions, predictability, tradeoff-surface shape) with a dollar estimate offered as an illustrative, clearly-caveated secondary figure rather than the headline metric.
